\subsection{Ewaluacja, ograniczenia i luki badawcze w automatycznym code review z użyciem LLM}
\label{subsec:llm_acr_evaluation_gaps}

Ocena tego, jak systemy ACR oparte na modelach językowych radzą sobie w praktyce, 
to w zasadzie jeden z najtrudniejszych tematów w literaturze. Wynika to z faktu, że code review to też proces silnie
społeczny. Tak jak już zauważono w sekcji \ref{subsec:llm_acr_prompting}, poleganie wyłącznie na miarach
ilościowych bywa mało miarodajne. Tufano i inni \cite{Tufano2024CodeReviewAutomation} trafnie punktują tutaj kwestię przygotowania zbiorów 
danych - jeśli mapowanie komentarzy do konkretnych zmian w kodzie jest nieprecyzyjne, to wyniki 
modelu mogą wyglądać na papierze bardzo dobrze, podczas gdy w realnym procesie pracy programisty będą 
zupełnie nieprzydatne. Prowadzi to do co raz częściej występujących głosów w literaturze, że konieczne jest używanie
również weryfikacji funkcjonalnej obok metryk automatycznych \cite{Tufano2024CodeReviewAutomation, Zhang2024SurveyLLMSE}.
Nie jest to jednak proste do wdrożenia w badaniach.

\paragraph{Metryki tekstowe a poprawność merytoryczna.}
Metryki generacji tekstu, takie jak BLEU czy BERTScore to kolejny aspekt warty omówienia. Haider 
i in. \cite{Haider2024PromptingFineTuningRCG} pokazali, że pozwalają one na pewną formę porównywania wyników między modelami, 
to tak naprawdę niewiele mówią one o tym, czy wygenerowana porada ma sens merytoryczny i czy jest zgodna z konkretnymi
standardami danego projektu \cite{Tufano2024CodeReviewAutomation}. Problem nie jest łatwy w rozwiązaniu, ponieważ
komentarze w procesie przeglądu kodu dotyczą skrajnie różnych rzeczy: czystości stylu, wydajność kodu, poprawność
logiki biznesowej i inne. Ciężko więc oczekiwać, że jakość recenzji da się 
sprowadzić do jakiegoś jednego, uniwersalnego wskaźnika liczbowego, co zresztą potwierdzają też autorzy pracy \cite{Li2025CodeDoctor}, 
wskazując na trudności w obiektywnej kategoryzacji takich uwag.

\paragraph{Ewaluacja funkcjonalna i ryzyko regresji.}
Praktyka bardziej skupiona na odpowiedzi, czy sugerowana zmiana faktycznie poprawia
kod, a nie tylko wygląda "rozsądnie" na papierze jest bliższe rozsądkowi. Protokół walidacji oparty o testy jednostkowe, zaproponowany
przez Cihana i in. i omówiony wcześniej przy okazji intencji modelu, pokazuje przy tym dość jasno, że ryzyko wprowadzenia
zmian szkodliwych nie jest zerowe \cite{Cihan2025EvaluatingLLMsForCodeReview}. W kontekście projektowanego systemu
oznacza to konieczność dołożenia elementów zabezpieczających, np. automatycznej walidacji w procesach CI
oraz podejścia \emph{human-in-the-loop}.

\paragraph{Ewaluacja diagnostyczna i znaczenie kontekstu.}
Faktycznie zrozumienie zmian w code review jest nowoczesnym podejściem które co raz częściej widać w systemach automatyzacji
tego procesu. Służą do tego dedykowane benchmarki, takie jak CodeReviewQA 
\cite{Lin2025CodeReviewQA}, które pozwalają sprawdzić, czy model w ogóle potrafi poprawnie zinterpretować intencje stojące za 
zmianami w kodzie, zanim jeszcze zaproponuje konkretne poprawki \cite{Lin2025CodeReviewQA}. Odpowiednie osadzenie modelu
w kontekście konkretnego repozytorium do kolejny istotny aspekt widoczny w literaturze. Zastosowanie mechanizmów RAG albo
łączenie LLM-ów z analizą statyczną — co opisują Meng \cite{Meng2025RARe} oraz Jaoua \cite{Jaoua2025LLMStaticAnalyzersCR} — nie tylko
poprawia samą jakość generowanych odpowiedzi, ale też sprawia, że łatwiej jest prześledzić ich pochodzenie.

\paragraph{Luki w badaniach i wnioski dla projektowania systemów.}
Przy analizowaniu obecnych rozwiązań ACR widać jednak kilka wyraźnych braków:
\begin{itemize}
    \item \textbf{Niepełny cykl pracy}: Większość prac skupia się właściwie tylko na tym, 
jak wygenerować komentarz do kodu. Jak zauważa Germano \cite{Germano2025VulnSLR}, brakuje tutaj szerszego spojrzenia, które 
obejmowałoby też etap uzasadniania decyzji czy późniejszą, automatyczną naprawę wykrytego problemu.
    \item \textbf{Problemy z jakością danych treningowych}: Same dane bywają problematyczne. Yu i 
in. \cite{Yu2025Desiview} słusznie punktują, że bez rygorystycznego filtrowania (np. poprzez destylację DRC), modele po 
prostu uczą się od ludzi mało efektywnych wzorców komunikacji, co potem przekłada się na 
słabą jakość sugestii.
    \item \textbf{Praktyczne bariery przy wdrożeniach}: Na koniec zostają kwestie czysto techniczne i 
biznesowe. Doświadczenia opisane przez Ramesha na przykładzie Ericssona \cite{Ramesh2025EricssonExperience} pokazują, że finalnie to nie 
parametry modelu, a koszty, czas odpowiedzi i trudność zintegrowania narzędzia z obecnym workflow programistów 
decydują o tym, czy system się przyjmie.
\end{itemize}

Biorąc powyższe pod uwagę, stwierdzić można, że budowa efektywnego narzędzia do wspomagania przeglądu kodu
wymaga więcej niż prostego generowania kodu. Kluczowe okazuje się tutaj 
połączenie selektywnego podejścia do kontekstu z bardzo dokładną, wielowarstwową 
procedurą ewaluacji. Ocena nie może ograniczać się tylko do samej poprawności składni; 
musi brać pod uwagę też realną użyteczność i pilnować, żeby nie dochodziło 
do regresji w kodzie \cite{Tufano2024CodeReviewAutomation, Zhang2024SurveyLLMSE}. 
Wspomniane aspekty są kluczowe dla zapewnienia stabilności narzędzia.